\documentclass[dvipsnames]{beamer}
\input{shared.tex}

%\usepackage{orcidlink}
%\usepackage{authblk}
\usepackage{fontawesome}

\title[Multilingual Metonymy \& Metaphor]{Metonymy is more multilingual than metaphor: \\ Analysing tropes using ChainNet and the Open Multilingual Wordnet}


\usepackage[normalem]{ulem}
\newcommand{\ul}{\uline}
\newcommand{\email}[1]{\href{mailto:#1}{#1}}
\author[Bond \& Maudslay]{Francis \ul{Bond}$^\spadesuit$ and Rowan Hall \ul{Maudslay}$^\diamondsuit$ } 

\institute[UPOL \& Cambridge]{$^\spadesuit$Department of Asian Studies and Sinofon Project\\
  Palacký University Olomouc \\[1ex]
 $^\diamondsuit$   Dept.\ of Computer Science \& Technology \\
  University of Cambridge \\[1ex]
  \email{bond@ieee.org rh635@cam.ac.uk}  
}
\date[2025-01-30 (GWC 2025)]{2025-01-30 \\
13th Global Wordnet Conference, Pavia}

% \usetheme{Boadilla}
% \usecolortheme{seahorse}%spruce}
% \usepackage{polyglossia}
% \setmainlanguage{english}
% \setsansfont{Linux Libertine O}
% \usepackage{xeCJK}
% \setCJKmainfont{Noto Sans CJK JP}
% \setCJKsansfont{Noto Sans CJK JP}
% \usefonttheme{professionalfonts}

% \usepackage[round]{natbib}
% \newcommand{\citeplanguageresource}{\citep}

\usepackage{expex}
\usepackage{makecell}
\newfontfamily{\ipafont}{Gentium}
\newcommand{\ipa}[1]{{\ipafont\selectfont #1}}
\usepackage{tabularx}

\usepackage[e,f,j]{mtg2e}
\newcommand{\cmn}{\mtciteform}
\newcommand{\kor}{\mtciteform}
\renewcommand{\mtcitestyle}[1]{\textcolor{teal}{\textit{#1}}}
\usepackage{xcolor}
 \newcommand{\bs}{\textbackslash}   % backslash
 \newcommand{\cmd}[1]{{\bf \color{red}#1}}   % highlights command
 \newcommand{\wemp}[1]{\ul{#1}}   % highlight
 \newcommand{\emp}[1]{\textbf{\color{blue}{#1}}}   % highlight
 \newcommand{\rel}[1]{\textsc{\color{blue}{#1}}}
 \newcommand{\lnk}[1]{\rel{#1}}
 \newcommand{\dff}[1]{\texttt{\color{blue}{#1}}}
  \newcommand{\ghost}[1]{\color{gray}{#1}} 
%\newcommand{\txx}[1]{\textbf{\textcolor{blue}{#1}}}

\usepackage{chainnet}
\usepackage{booktabs}
\newcommand{\clex}[1]{\textbf{\texttt{#1}}}
\newcommand{\ldiff}[1]{\textbf{\texttt{#1}}}
\newcommand{\dom}[1]{\textbf{\texttt{#1}}}

\newcommand{\LEX}[6][]{\textbf{\textit{#2}} \textit{#1} (#3: #4) #5 --- \textit{#6}}

\begin{document}

\maketitle

\begin{frame}{Metaphor and Metonymy}

  \begin{itemize}
  \item Metonymy and metaphor are two fundamental tropes (figures of speech)
  \item Both play crucial roles in the way we understand and use language
  \item \txx{Metonymy} involves the substitution of one term for
    another to which it is closely related
    \begin{itemize}
    \item \eng{I ate a \ul{chestnut}}
    \item \eng{I pruned the \ul{chestnut} [tree]}
    \item \eng{This table is made from \ul{chestnut} [wood]}
    \item Metonymy works by contiguity or association between concepts,
      usually by some shared attribute, or by a part-whole relation
    \end{itemize}
  \item \txx{Metaphor} involves understanding one concept in terms of
    another, often unrelated, concept.
     \begin{itemize}
    \item \eng{This horse has a \ul{chestnut} on its leg} ``a small horny bump''
    \item Metaphors create new concepts in the target domain by
      mapping elements from the source domain
    \end{itemize}
    \item They are crucial in understanding novel uses of words
  \end{itemize}

\end{frame}




\begin{frame}{ChainNet}

  \begin{itemize}
  \item ChainNet is an attempt to comprehensively model these tropes \citep{maudslay-etal-2024-chainnet-structured,Bond:Maudslay:2025}
    \begin{itemize}
    \item All nouns (word + synset combinations) in the
      Princeton Wordnet \citep{_Fellbaum:1998} with three or more senses$^*$
      were annotated
  \item Every sense linked to another sense, or treated as a homonym (unrelated) 
  \item Most annotated by 1 annotator, some by 2 annotators
    \begin{itemize}
    \item Inter-annotator agreement 70\% (88\% for same prototype)
    \item Intra-annotator agreement 81\% (92\% for same prototype)
    \end{itemize}
  \item Originally grew out of work in determining homonyms --- which senses of a word are related
  \end{itemize}
  \item There was some early work on metaphor$^{**}$  \citep{alonge-2006-italian} and more recent work on metonymy \citep{khishigsuren_bella_brochhagen_marav_giunchiglia_batsuren_2022} which we discuss briefly  in the paper.

\end{itemize}

\begin{small}
 $^*$ almost, and some with two senses.   All should eventually be annotated \\
$^{**}$ if anyone knows where the Italian Metaphor Database is, please let us know
\end{small}
\end{frame}


\begin{frame}{chestnut}
    \begin{center}
\begin{tikzpicture}[derivation]

\litbox[4cm]{1}{at (0,0)}{\sense{chestnut}{3} edible nut of any of various chestnut trees}

\assbox[4cm]{3}{[below right = -5mm and 2cm of 1]}{%
\sense{chestnut}{2} any of several attractive deciduous trees yellow-brown in autumn; yield a hard wood and edible nuts in a prickly bur}

\assbox[2cm]{4}{[below right = 0.85cm and -1cm of 1]}%
{\sense{chestnut}{4} the brown color of chestnuts}

\assbox[3cm]{6}{[below  = 8mm of 4]}{\sense{chestnut}{6} a dark golden-brown or reddish-brown horse }

\metbox[3cm]{5}{[below left = 0.8cm and -3cm of 1]}%
{\sense{chestnut}{5} a small horny callus on the inner surface of a horse's leg}

\assbox[4cm]{2}{[below = 13mm  of 3]}{\sense{chestnut}{1}  [\synonym{chestnut tree}] wood of any of various chestnut trees of the genus Castanea }

   \path[related] (1) edge node[anchor=south, font={\footnotesize}] {\relatededgelabel} (3);

%  \node[shape=coordinate] (origin) [above = .8cm of 4, node contents={}];

  \path[related] (3) edge node[anchor=west, font={\footnotesize}] {\relatededgelabel} (2);

    % \node[shape=coordinate] (origin2) [above left = .8cm and -.5cm of 2, node contents={}];
    % \node[shape=coordinate] (origin3) [below = .8cm of origin2, node contents={}];

 % \path[analogy] (origin2) edge node[anchor=west, font={\footnotesize}] {\metaphoricaledgelabel} (origin3);

 \path[analogy] (1) edge node[anchor=south, xshift = -8mm, yshift=-3mm, font={\footnotesize}] {\metaphoricaledgelabel} (5);
 %\node[shape=coordinate] (origin) [above = .4cm of 5, node contents={}];
 
 \path[related] (1) edge node[anchor=south,  xshift = 9mm, yshift=-3mm, font={\footnotesize}] {\relatededgelabel} (4);

 \path[related] (4) edge node[anchor=south,  xshift = 9mm, yshift=-3mm, font={\footnotesize}] {\relatededgelabel} (6);

%  \node[shape=coordinate] (origin) [above = .4cm of 5, node contents={}];

\end{tikzpicture}
\end{center}

\end{frame}



\begin{frame}{Can we see how English metaphors are used cross-lingually?}
  \begin{itemize}
  \item For example, are translations of linked pairs likely to be
    translated with the same word
    \begin{itemize}
    \item If so then metaphor/metonomy holds cross language
    \item[\ldots] or this may indicate issues with automatic construction of wordnets
    \end{itemize}
  \item For example \eng{head} ``body part'' is metaphorically
    extended to ``person in charge'' in English and this is also true in Japanese,
    (頭) Italian (\eng{capo}) and many other languages
  \item But not all metaphors are shared:  \eng{head} ``main part of a grammatical constituent'' is not 頭 in Japanese
    \bigskip
  \item ChainNet is made for English, but we can expect many of the
    tropes to work with other languages
  \item And we have many wordnets linked to English in the Open Multilingual Wordnet
    

  \end{itemize}
\end{frame}

\begin{frame}{Measure the translation overlap with OMW}

  \begin{itemize}
  \item For every English word that has been annotated
    \begin{itemize}
    \item For each pair of English senses
    \item Look up all translations of each sense
      \\ measure the overlap (we use Jaccard similarity)
      \item Store and note if the senses are linked or not
    \end{itemize}
  \item So for example \eng{head} has 3 senses (really many more)
    \begin{enumerate}
    \item head, caput
      \\ ヘッド, \ul{頭}, 頭部 
    \item head, chief, top dog
      \\ 大頭, 主任者, 御頭, 頭領, \ul{頭}
    \item head,  head word
      \\ 主要語
    \end{enumerate}
  \item Overlap
    \begin{itemize}
    \item 1-2: 0.125 %(/ 1.0  8) 0.125
    \item 2-3: 0
    \item 1-3: 0
    \end{itemize}
  \end{itemize}
  

\end{frame}

\begin{frame}[allowframebreaks]{We do this for all the lexicons in OMW-1.4}

 
  \begin{center}
  \begin{tabular}{lrrrr}
    Lang & Unlinked & Metaphor & Metonomy & Non-Zero Instances \\ \hline
sv & 0.70 & 1.71 & 2.10 & 417   \\
he & 0.69 & 1.58 & 2.31 & 472   \\
eu & 0.82 & 1.30 & 1.80 & 5,562  \\
ja & 0.89 & 1.11 & 1.55 & 8,130 \\
ro & 0.83 & 1.30 & 1.71 & 6,429 \\
sq & 0.74 & 1.27 & 2.38 & 822 \\
%en & 0.97 & 1.08 & 1.07 & 33415 \\
da & 0.74 & 1.46 & 2.09 & 332 \\
lt & 0.63 & 1.46 & 2.86 & 692 \\
 \end{tabular}
\end{center}

\begin{itemize}
\item \txx{unlinked} has no link
\item \txx{metaphor} is linked by metaphor
\item \txx{metonymy} is linked by metonymy
\end{itemize}
Normalize by dividing  by the average distance for all senses.

  \begin{center}
  \begin{tabular}{lrrrr}
    Lang & Unlinked & Metaphor & Metonomy & Non-Zero Instances \\ \hline
es & 0.85 & 1.17 & 1.79 & 5,915 \\
 hr & 0.81 & 1.22 & 1.98 & 3,576 \\
    fr & 0.94 & 1.10 & 1.25 & 18,975 \\
nb & 0.73 & 1.49 & 2.12 & 339 \\
bg & 0.73 & 1.54 & 2.11 & 380 \\
gl & 0.70 & 2.13 & 1.55 & 271 \\
cmn & 0.74 & 1.57 & 1.97 & 1,347 \\
it & 0.80 & 1.39 & 1.81 & 4,614 \\
arb & 0.74 & 1.28 & 2.38 & 1,337 \\
pl & 0.75 & 1.43 & 2.07 & 1,326 \\
el & 0.72 & 1.30 & 2.45 & 1,274 \\
iwn & 0.79 & 1.46 & 1.80 & 968 \\
nn & 0.73 & 1.52 & 2.10 & 340 \\
id & 0.87 & 1.21 & 1.58 & 9,803 \\
  \end{tabular}
\end{center}

    \begin{center}
      \begin{tabular}{lrrrr}
    Lang & Unlinked & Metaphor & Metonomy & Non-Zero Instances \\ \hline
is & 0.73 & 1.51 & 2.13 & 388 \\
ca & 0.83 & 1.23 & 1.84 & 6,000 \\
sk & 0.73 & 1.35 & 2.31 & 2,067 \\
fi & 0.82 & 1.35 & 1.76 & 7,971 \\
nl & 0.79 & 1.38 & 1.90 & 2,991 \\
sl & 0.81 & 1.41 & 1.74 & 6,243 \\
th & 0.80 & 1.48 & 1.71 & 2,336 \\
pt & 0.80 & 1.26 & 1.96 & 5,414 \\
        zsm & 0.87 & 1.22 & 1.60 & 10,272 \\
        \midrule
Mean & 0.78 & 1.39 & 1.96 &  3,774\makebox[0mm]{~~~.3} 
  \end{tabular}
    \end{center}
    
  Unlinked  share fewer translations, metaphor shares more and metonomy shares much more.

  \Large Metonomy is more likely to hold cross linguistically!
  
\end{frame}

\begin{frame}{Discussion}

  \begin{itemize}
  \item The wordnets are of vastly different sizes: the number of sense
    pairs that have a translation varies from 271 (Galician) to 18,975 (French). 
  \item   Only Galician has the     score for metaphor (2.13) larger the score for metonymy (1.55), probably due to data sparsity. We would expect it to behave much like
    Portuguese (1.26/1.96).
  \item In order to measure perfectly how well tropes carry
    over between languages, we would need to mark metonymy and metaphor
    systematically for each language, and make sure all synsets for all
    senses have all relevant translations, \ldots 
  \item Without exception, senses linked by tropes are more likely to have an identical translation that those senses which are not linked
  \item With one exception, metonymy is more likely to have  an identical translation then metaphor  
  \end{itemize}

\end{frame}

\begin{frame}{Translations of the senses of \protect\eng{cherry}}
  \begin{center}
 \small 
  \begin{tabular}{llllll}% \toprule
    Sense & Mandarin & Japanese & Finnish & Italian & Portuguese \\ \midrule

    \sense{cherry}{1}   & 樱桃木 & 桜,   \ul{桜材} & kirsikkapuumetsä & ciliegio  & cerejeira \\ 

    \sense{cherry}{2}  & 樱桃树 & 桜, 櫻 & kirsikkapuu & ciliegio  &  cerejeira\\

    \sense{cherry}{3}  & 樱桃  & 桜桃, 桜ん坊 & kirsikka & cerasa , ciliegia  & ginja , cereja \\ 
    
    \sense{cherry}{4} &  \ul{樱桃红} & \ul{桜ん坊色}   &  kirsikanpunainen & ciliegia,  \ghost{rosso ciliegia}\ & cereja\ %\bottomrule
  \end{tabular}
\end{center}

\begin{itemize}
\item When we look at the translations, we see that the metonymy that is unmarked in English, is marked in other languages
\item Even within English, other synonyms have information
\end{itemize}

\end{frame}


\begin{frame}{Can we use other lemmas to classify links}

  \begin{itemize}
  \item Synsets have sets of synonyms, so the other senses may include useful information
  \item Consider the synsets with a sense of chestnut
    \begin{itemize}
    \item[s1] chestnut ``edible nut of any of various chestnut trees of the genus Castanea'' (NATURAL\_BODY)

    \item[s2] chestnut, chestnut tree ``any of several attractive deciduous trees yellow-brown in autumn; yield a hard wood and edible nuts in a prickly bur'' (PLANT)
    \item[s3] chestnut ``wood of any of various chestnut trees of the genus Castanea''
      (CHEMICAL!)
    \end{itemize}
  \item These are annotated as s1 → s2 → s3 (metonomy)
  \item If we look at all the senses
    \begin{itemize}
    \item  s1 → s2 has a diff: ``+ tree''
    \item s2 → s3 has a diff: ``- tree''
    \end{itemize}
  \item These patterns repeat \ldots
    \\ Cherry → Cherry Tree → Cherry\hfill (en)
  \end{itemize}

 
\end{frame}
\begin{frame}{They can be found across languages}

  \begin{itemize}
  \item s1 → s2: kastanje  →  kastanjetræ \hfill (da)
  \item s1 → s2: chestnut  →  chestnut tree \hfill (en)
  \item s1 → s2: marron →	marronnier \hfill (fr)
  \item s1 → s2: kastanja →	kastanjapuu \hfill (fi)
  \item  s1 → s2:	kastanje →	kastanjeboom \hfill (nl)
  \item  s1 → s2:	castană →	castan \hfill (ro)
  \item  s1 → s2:	castanăcomestibilă →	castan \hfill (ro)
    
  \end{itemize}	
\end{frame}

% \begin{frame}{Different languages mark different things}
%     \begin{itemize}
%     \item 
%       accounting, accountancy
%       ``the occupation of maintaining and auditing records and preparing financial reports for a business''
% \\      アカウンティング, 経理, \ul{会計}, 会計業務, 主計 
%     \item accounting 
%       ``a system that provides quantitative information about finances''
% \\      アカウンティング, 経理, 会計, 簿記, 主計, \ul{会計学}, 計数 
% \item history ``all that is remembered of the past as preserved in writing; a body of knowledge''
% \ul{歴史}, 青史, ヒストリ, 史, 史乗
    	
%       \item history ``the discipline that records and interprets past events involving human beings''
% 歴史, 史学, \ul{歴史学}  
  
% \end{itemize}

% Japanese marks fields of study with   \ul{学}.   So we can use it to sub-type the unmarked metaphor relation in English, \ldots
% \end{frame}

% \begin{frame}{Issues}
  
% \begin{itemize}
% \item Some of the derivation links are spelling variants
%   \\ ('i60334', 'toilet', 'i39761', 'toilette'), ('i73661', 'duet', 'i80460', 'duette'), ('i109102', 'quartet', 'i79264', 'quartette'), ('i80466', 'sextet', 'i73665', 'sextette'), \ldots
%   \begin{itemize}
%   \item  In OMW 2.0 these should not be separate senses (just orthographic variants)
%   \item Someone has to discover and fix the entries
%   \end{itemize}
% \item There are many small sub-regularities, it would be good to capture cheaply
%   \\	('i109106', 'eight', 'i53316', 'eight-spot'), ('i109110', 'ten', 'i60122', 'ten-spot'), ('i109105', 'seven', 'i58705', 'seven-spot'), ('i109102', 'four', 'i54015', 'four-spot'), ('i109104', 'six', 'i58979', 'six-spot'), ('i109107', 'nine', 'i56605', 'nine-spot'), ('i109103', 'five', 'i53789', 'five-spot')]
%  \\ \textbf{metonomy:-spot}
%   \\ and why are we missing two-spot ``deuce'', three-spot?
% \end{itemize}

% \end{frame}
\begin{frame}{Ongoing Research to Find Patterns}
  
\begin{itemize}
\item We have made a searchable database where the edges are decorated by the morphological differences
  \begin{itemize}
  \item So we can see all pairs linked by \dff{en:+ tree} or \dff{da:+træ} 
  \item The nodes link to synsets, so we can look
    up supertypes: we can say this link is of type  \lnk{food:plant}
  \\ or  \lnk{plant is food} in the cognitive linguistic style
  \item In this way we can group the tropes together in families
  \item Similar to \citet{Stoyanova:Mititelu:Leseva:Iordachioaia:2025}
  \end{itemize}
\item Other people have done this top down \cite[e.g start at \lnk{food:plant} and looked for examples]{khishigsuren_bella_brochhagen_marav_giunchiglia_batsuren_2022} 
\item With ChainNet we can go bottom-up and find many more patterns
  \begin{itemize}
  \item We can use multi-lingual features
  \item We need to investigate methods of clustering
  \item We need to find a good way of evaluating
  \end{itemize}
\item We would like to also map to existing Conceptual Metaphors
\end{itemize}

\end{frame}

\begin{frame}{The Patterns can be used Generatively}
  \begin{itemize}
  \item If we see a word that is of type \lnk{food}, then we can expect that
    there may be a kind of plant using the same word
    \begin{itemize}
    \item Useful for analyzing texts that otherwise don't make sense
    \item Can be used to analyse humour and word play \\
    \end{itemize}
  \end{itemize}

  \begin{center}
    \includegraphics[width=0.5\textwidth]{spaghetti-tree.png} \\
    \href{https://en.wikipedia.org/wiki/Spaghetti-tree_hoax}{BBC Spaghetti Tree hoax (1957)}
  \end{center}
\end{frame}

\begin{frame}{Metaphor in Large Language Models}
  \begin{itemize}
  \item Because these tropes are so established, LLMs capture the
    patterns well, and can generalise to novel examples, at least for languages with a lot of online data
  \item E.g. For Spanish, \citet{puraivan2024metaphor} found ChatGPT
    could differentiate literal from metaphorical senses with an
    accuracy of 85-88\% (on a small dataset)
    \begin{itemize}
    \item This is high enough to cause \txx{automation blindness}
    \item It is very hard for a human to spot mistakes
    \end{itemize}
  \item We expect that this accuracy will go down for languages with
    less data and for less common senses, but have not yet confirmed
    this
  \item There are also interesting questions of how much metaphor
    interpretation in multilingual models is effected by other languages
    \\ The DanNet group have found metaphors in Danish not ofund in English are less likely to be recognized
  \item So we feel it is important to
    keep working on producing reliable
    manually annotated data
  \end{itemize}%     

  
 \end{frame}


\begin{frame}{Future Work}

  \begin{itemize}
  \item We are interested in looking at
    the relation between in-context uses
    and tropes
  \\  Using sense-tagged corpora we can see 
    \begin{itemize}
    \item If basic uses or extended uses are more common
    \item Do similar tropes cluster?
    \item Do different genres use different tropes?
    \end{itemize}
  \item We would like to do large scale
    annotation on another language
    (probably start with Japanese)
  \item We are looking in detail at a small sample: currently investigating
    \lnk{animal:person} metaphors
    \begin{itemize}
    \item  \eng{Be a \ul{lamb} and give me a hand, ...}
    \item In Chinese, Czech, English, French,
      Indonesian, Japanese, Javanese,
      Korean, Slovak and Vietnamese  --- it would be good to include some sign
    \end{itemize}
  \item We would like to investigate similarities and differences in images, \ldots
  \end{itemize}
\end{frame}


\begin{frame}{Summary}
  
\begin{itemize}
\item The combination of detailed annotation of meaning and
  cross-lingual links allow us to explore differences in meaning
  representation
  \begin{itemize}
  \item  Metonymy is based on association and contiguity
    \\ \txx{more universal}
  \item  Metaphor often requires analogical reasoning and cultural framing
    \\ \txx{more culture and language specific}
  \end{itemize}
\item We don't know if the tropes are independently discovered, come from a common ancestor, or are borrowed (the same problem as with cognates in etymology)
\item We need to add to wordnet and make searchable
  \begin{itemize}
  \item Add as sense-level relations, not etymology, contra \citet{khan-etal-2020-modelling}
  \end{itemize}
\item LLMs may capture these regularities, but are hard to explore with
\end{itemize}

\end{frame}

\begin{frame}{Acknowledgements}
  
\begin{itemize}
\item Thanks to Bolette Pederson and the DanNet group for some very useful discussions
\item ChainNet is primarily the work of Rowan Hall-Maudslay \email{r.hallmaudslay@gmail.com}
  \\ \url{https://github.com/rowanhm/ChainNet}
\item The Open Multilingual Wordnet brings together wordnets made by many projects, we thank them all
\item Data and code for the cross-lingual analysis is available at \url{https://github.com/bond-lab/chainnet-xling}
  \\ (still changing a lot)
\end{itemize}

\end{frame}

\begin{frame}{OMW 1.4}
  \begin{itemize}
  \item The OMW is a collection of wordnets for different languages.
  \item These languages are linked through the collaborative interlingual index \citep[CILI;][]{Bond:Vossen:McCrae:Fellbaum:2016}, a language-neutral list of concepts that is linked to by the various wordnets.
    
  \item The combined wordnets include 
English \citep{_Fellbaum:1998}, 
Albanian \citep{Ruci:2008}, 
Arabic \citep{Black:Elkateb:Rodriguez:Alkhalifa:Vossen:Pease:Bertran:Fellbaum:2006}, 
Chinese \citep{Huang:Hsieh:Hong:Chen:Su:Chen:Huang:2010}, 
Danish \citep{Pedersen:2009}, 
Finnish \citep{Linden:Carlson:2010}, 
French \citep{Sagot:Fiser:2008}, 
Hebrew \citep{Ordan:Wintner:2007},
Indonesian and Malaysian (Nurril Hirfana et al., \citeyear{Noor:Sapuan:Bond:2011}), 
Italian \citep{Pianta:Bentivogli:Girardi:2002,Toral:Bracale:Monachini:Soria:2010}, 
Japanese \citep{Isahara:Bond:Uchimoto:Utiyama:Kanzaki:2008}, 
Norwegian Bokm\r{a}l and Norwegian Nynorsk (Lars Nygaard, personal communication 2012),
Persian \citep{Montazery:Faili:2010}, 
Portuguese \citep{Paiva:Rademaker:2012};
Polish \citep{Piasecki:Szpakowicz:Broda:2009}, 
Thai \citep{Thoongsup:Charoenporn:Robkop:Sinthurahat:Mokarat:Sornlertlamvanich:Isahara:2009},
and Basque, Catalan, Galician and Spanish \citep{Gonzalez-Agirre:Laparra:Rigau:2012}.
\item Data is available from \url{https://github.com/omwn/omw-data}
\item We used the python \defn{wn} module \citep{Goodman:Bond:2021}
\end{itemize}
\end{frame}


\begin{frame}[allowframebreaks]{References}
  \bibliographystyle{acl_natbib}
  \bibliography{mtg.bib,local.bib}
\end{frame}

\end{document}

%%% Local Variables:
%%% coding: utf-8
%%% mode: latex
%%% TeX-PDF-mode: t
%%% TeX-engine: luatex
%%% TeX-master: t
%%% End: 

